\documentclass[11pt]{article}
\usepackage[margin=1in]{geometry}
\usepackage{amsmath,amssymb,amsthm}
\usepackage[T1]{fontenc}
\usepackage{lmodern}
\usepackage[hidelinks]{hyperref}
\newtheorem{theorem}{Theorem}
\newtheorem{lemma}[theorem]{Lemma}
\newtheorem{proposition}[theorem]{Proposition}
\newcommand{\D}{\mathbb D}
\newcommand{\SR}{S_{\mathbb R}}
\newcommand{\cl}[1]{\overline{#1}}
\DeclareMathOperator{\Impart}{Im}
\setlength{\parskip}{0.25em}
\title{A qualitative pre-Schwarzian improvement\protect\\ for real schlicht functions}
\author{Alec Kriebel\\\small Independent researcher\\
\small\href{https://orcid.org/0009-0001-9320-500X}{ORCID: 0009-0001-9320-500X}}
\date{October 15, 2026}
\begin{document}
\maketitle
\begin{abstract}
For normalized univalent functions with real Taylor coefficients, the classical
centered pre-Schwarzian disk can be strictly shrunk at every fixed nonreal point,
uniformly over the class. On the real diameter its classical radius remains
sharp. We prove this qualitative answer to Hayman--Lingham Problem~6.86 by
combining classical coefficient equality, reflection symmetry of the equality
domains, and compactness. The gap is uniform on compact nonreal sets. No explicit
upper-gap formula or sharp radius is determined. The proof uses classical
methods, and no historical firstness is claimed.
\end{abstract}

\section{The fixed-point question}
Let $S$ be the holomorphic injections $f:\D\to\mathbb C$, where
$\D=\{z:|z|<1\}$, normalized by $f(0)=0$ and $f'(0)=1$.
Let $\SR$ be the subfamily with real Taylor coefficients, equivalently
$f(\cl z)=\cl{f(z)}$. Its members have nonvanishing derivative. Write
\[
 P_f=\frac{f''}{f'},\qquad
 A_f(z)=(1-|z|^2)P_f(z)-2\cl z.
\]
The classical bound is $|A_f(z)|\le4$ for $f\in S$.
Problem~6.86 of Hayman and Lingham \cite[printed pp.~114 and 147]{HL}
asks whether the corresponding estimate
\begin{equation}\label{source}
 \left|zP_f(z)-\frac{2|z|^2}{1-|z|^2}\right|
 \le \frac{4|z|}{1-|z|^2}
\end{equation}
improves when the coefficients are real. The source fixes the evaluation point;
its wording does not explicitly require a sharp or numerical formula.
We answer that qualitative question, with the necessary real-axis exception.
Postcomposition by a real affine map normalizes any real-on-real holomorphic
injection without changing $P_f$.

\begin{samepage}
\begin{theorem}\label{main}
For every $z\in\D$ the maximum
\[
 M(z)=\max_{f\in\SR}|A_f(z)|
\]
exists. It satisfies
\[
 M(z)<4\quad(\Impart z\ne0),\qquad M(x)=4\quad(-1<x<1).
\]
For each nonempty compact $K\subset\D\setminus\mathbb R$ there is
$\epsilon_K>0$ such that $|A_f(z)|\le4-\epsilon_K$ for all
$f\in\SR$ and $z\in K$.
\end{theorem}
\end{samepage}

Thus, for fixed nonreal $z$, the right side of \eqref{source} may be replaced by
$(4-\epsilon(z))|z|/(1-|z|^2)$, where $\epsilon(z)=4-M(z)>0$.
This is a uniform-in-function existence statement, not an evaluation of $M$.
At nonzero real points the classical radius is sharp even if the disk's center
is allowed to move: the examples below give its two opposite endpoints.
At $z=0$ the expression in \eqref{source} is identically zero; the auxiliary
quantity $A_f(0)=f''(0)$ is different.

\section{Classical equality and real symmetry}
We include the coefficient argument so that the equality information used in
the proof is explicit. The methods and classical estimate are not claimed as
new; the source refers to Hayman's classical treatment \cite{Hayman}.

\begin{lemma}\label{coefficient}
If $F(z)=z+a_2z^2+\cdots\in S$, then $|a_2|\le2$.
Equality holds exactly for $F(z)=z/(1-\lambda z)^2$, $|\lambda|=1$.
\end{lemma}
\begin{proof}
The zero-free function $F(z)/z$, extended at zero by $1$, has a logarithm $L$
on $\D$ with $L(0)=0$. On $|w|>1$ put
\[
 H(w)=w\exp[-L(w^{-2})/2]=w-\frac{a_2}{2w}+O(w^{-3}).
\]
It is odd and nonzero. If $H(w_1)=H(w_2)$, then squaring and using injectivity
of $F$ gives $w_2=\pm w_1$; oddness and nonvanishing exclude the minus sign.
Hence $H$ is an exterior injection.
For an exterior injection $w+b_0+\sum_{n\ge1}b_nw^{-n}$, Green's formula on
the image of $|w|=R>1$ gives the enclosed area
\[
 \pi\left(R^2-\sum_{n\ge1}n|b_n|^2R^{-2n}\right)\ge0.
\]
The curve is positively oriented since its exterior maps to the unbounded
component. Laurent series and their derivatives converge uniformly on that
circle. Each finite partial sum of this nonnegative series is at most $R^2$.
Letting $R\downarrow1$ and then increasing the number of terms gives
$\sum n|b_n|^2\le1$.
Here $b_1=-a_2/2$, proving the bound. Equality forces every further negative
coefficient to vanish; oddness removes $b_0$. Thus $H(w)=w-\lambda/w$,
$|\lambda|=1$, and squaring gives the stated $F$.
Conversely these maps are injective: equality of two values factors as
$(z-w)(1-\lambda^2zw)=0$. They attain $|a_2|=2$.
\end{proof}

For $a\in\D$ the normalized Koebe transform
\[
 T_af(w)=\frac{f((a+w)/(1+\cl a w))-f(a)}{(1-|a|^2)f'(a)}
\]
belongs to $S$, and $2a_2(T_af)=A_f(a)$. Lemma~\ref{coefficient} proves
$|A_f(a)|\le4$. If equality holds, $T_af$ is a Koebe rotation, so
$f(\D)$ is a complex affine image of a one-ray complement. This classifies
the image of the original $f$; it does not directly identify $f$ as a Koebe
rotation. For example the identity
\[
 \frac{z}{(1-z)^2}
 =\frac14\left[\left(\frac{1+z}{1-z}\right)^2-1\right]
\]
shows that its image omits exactly $(-\infty,-1/4]$; rotation gives the
general one-ray image.

Suppose now $f\in\SR$ and $|A_f(a)|=4$. Reflection symmetry makes its omitted
ray invariant under conjugation. The ray's unique finite endpoint is real,
and its direction must also be real. Since zero belongs to the image, the
ray is either $(-\infty,c]$, $c<0$, or $[c,\infty)$, $c>0$.
Put
\[
 k_+(z)=\frac{z}{(1-z)^2},\qquad k_-(z)=\frac{z}{(1+z)^2}.
\]
The appropriate map $g=4|c|k_+$ or $4|c|k_-$ has the same image as $f$ and
fixes zero. The inverse composition $g^{-1}\circ f$ is a disk automorphism
fixing zero; its derivative has modulus $1/(4|c|)=1$. Hence $|c|=1/4$,
and the positive normalized derivatives force $f=k_+$ or $f=k_-$.
Direct differentiation yields
\begin{equation}\label{deficit}
 |A_{k_+}(z)|^2=|A_{k_-}(z)|^2
 =16-\frac{48(\Impart z)^2}{|1-z^2|^2}.
\end{equation}
Thus equality is impossible at a nonreal point in $\SR$. On the real diameter,
$A_{k_+}(x)=4$ and $A_{k_-}(x)=-4$.

\begin{proof}[Proof of Theorem~\ref{main}]
We still need to pass from separate strictness for every $f$ to a strict
maximum over all $f$. The universal bound gives
\[
 |P_f(re^{i\theta})|\le\frac{4+2r}{1-r^2}.
\]
Integrating the radial derivative of $\log|f'|$, and then $f'$, gives
\[
 |f'(z)|\le\frac{1+|z|}{(1-|z|)^3},\qquad
 |f(z)|\le\frac{|z|}{(1-|z|)^2}.
\]
Montel's theorem makes $S$ a normal family. A locally uniform limit retains
derivative $1$ at zero and is nonconstant. To check injectivity, fix $w\in\D$
and apply Hurwitz to the zero-free holomorphic divided differences
$(f_n(z)-f_n(w))/(z-w)$, extended at $w$ by $f_n'(w)$.
Their limit cannot be identically zero, since that would make the limit of
$f_n$ constant. Thus the limit remains injective, with nonzero derivative.
Consequently $S$ is compact, and reflection makes $\SR$ a closed compact
subfamily. Derivative convergence and a nonzero limiting derivative imply
joint continuity of $(f,z)\mapsto|A_f(z)|$ at interior points.
The fixed-point maximum is attained, and the equality analysis above excludes
$4$ at every nonreal point. On $\SR\times K$ joint continuity similarly
gives an attained maximum below $4$. The real Koebe maps prove sharpness on
the real diameter. (For empty $K$ the assertion is vacuous.)
\end{proof}

\section{The remaining radius and two boundary obstructions}
The following exact reductions describe what the qualitative result leaves
undetermined. They do not provide an upper-gap formula.

\begin{proposition}\label{reduction}
Let $\delta(z)=|\Impart z|/|1-z^2|$. The function $M(z)$ depends only on
$\delta(z)\in[0,1/2)$, and
\begin{equation}\label{bounds}
 \max\{4\sqrt{1-3\delta(z)^2},\,8\delta(z)\}\le M(z)<4
 \quad(\Impart z\ne0).
\end{equation}
There is no single positive gap valid at all nonreal points of $\D$.
\end{proposition}
\begin{proof}
For real $a\in(-1,1)$, write $\phi_a(w)=(w+a)/(1+aw)$.
The transform $T_a$ is a bijection of $\SR$, with inverse $T_{-a}$.
The chain rule gives
\[
 A_{T_af}(w)=\frac{1+a\cl w}{1+aw}\,A_f(\phi_a(w)).
\]
The phase has modulus one, so $M(w)=M(\phi_a(w))$. Also
$M(\cl w)=M(w)$. The same maps preserve $\delta$, and
\[
 |1-z^2|^2=(1-|z|^2)^2+4(\Impart z)^2.
\]
For $x=\operatorname{Re}z$, $r=|z|$, choose
\[
 a=\frac{2x}{1+r^2+\sqrt{(1+r^2)^2-4x^2}}.
\]
Its modulus is below one, and the real part of $(z-a)/(1-az)$ vanishes,
because $xa^2-(1+r^2)a+x=0$. Thus a nonreal $z$ is equivalent to $it$ or
$-it$, where $0<t<1$ and
\[
 \delta=\frac{t}{1+t^2},\qquad
 t=\frac{2\delta}{1+\sqrt{1-4\delta^2}}.
\]
This proves the one-variable reduction. Equation~\eqref{deficit} gives the
first lower bound. For $h(z)=z/(1-z^2)$, the factorization
\[
 h(z)-h(w)=\frac{(z-w)(1+zw)}{(1-z^2)(1-w^2)}
\]
proves univalence on $\D$, and $A_h(it)=8it/(1+t^2)$. The real transforms
give the second lower bound at every point. The upper bound is
Theorem~\ref{main}. As $\delta\to0$, the Koebe lower bound tends to $4$;
as $\delta\to1/2$, the other one tends to $4$, realized along $z=it\to i$.
Therefore no global positive gap exists.
\end{proof}

The lower-bound branches cross at $\delta=1/\sqrt7$.
At $z=i/2$ their squares are $208/25$ and $256/25$, respectively, so the two
real Koebe maps cannot alone be extremizers there. These are lower bounds,
not a proposed formula for $M$. The first limiting obstruction is
$\delta\to0$, for example on approaching an interior real point; merely
$\Impart z\to0$ near the boundary endpoints $\pm1$ does not suffice.

\section{Scope, prior work, and verification}
The result is a qualitative application of classical equality rigidity and
compactness to the fixed-point question. It determines neither an explicit
positive upper gap nor the sharp function $M(it)$, its general extremizers,
or the full variability region. The 2018 source's dated no-progress update
does not certify present openness.

\noindent\begin{minipage}{\linewidth}
Relevant earlier work includes Koepf's classification of extreme points
\cite[Section~4]{Koepf}; its linear-functional setting does not directly
optimize $f''/f'$. Rodemich \cite{Rodemich} announced an argument bound for
$(1-z^2)f'(z)$ in the same real class; differentiating such an inequality
does not establish the present derivative bound. The supplement records a
precise test of that proposed shortcut, not a counterexample within the globally
univalent class. Jenkins's own survey \cite[p.~7]{Jenkins1962} identifies derivative
bounds in \cite{Jenkins1960}. The full 1960 article and the solved cases of
\cite{AlePop} were not obtained in the bounded priority audit. Their exact
overlap remains unverified. Accordingly no historical firstness, new general
method, exhaustive literature clearance, or current-openness guarantee is
claimed. These limits do not enter the self-contained proof above.
\end{minipage}

The accompanying verification archive preserves the original research
candidate and its controlling scope clarification, independent analytic and
source audits, bounded priority records, and exact algebraic controls.
The analytic argument is a written proof; the controls are not a formal
proof certificate. Saved exploratory Loewner values are explicitly
uncertified and play no role in the proof. Full third-party source scans
are excluded from the archive.

\paragraph{AI assistance and review status.}
AI tools were used extensively in developing the argument, checking sources
and computations, adversarial verification, and writing this note and its
supporting materials. This is an unrefereed preprint. Internal AI reviews do
not constitute independent human peer review.

\begin{thebibliography}{9}
\bibitem{HL} W.~K. Hayman and E.~F. Lingham,
\emph{Research Problems in Function Theory (New Edition)},
\href{https://arxiv.org/abs/1809.07200v2}{arXiv:1809.07200v2} (2018),
Chapter~6, printed pp.~114 and 147, Problem~6.86.
\bibitem{Hayman} W.~K. Hayman, \emph{Multivalent Functions},
Cambridge Tracts in Mathematics and Mathematical Physics~48,
Cambridge University Press, 1958.
\bibitem{Koepf} W. Koepf, On nonvanishing univalent functions with real
coefficients, \emph{Math. Z.} \textbf{192} (1986), 575--579.
\href{https://doi.org/10.1007/BF01162704}{doi:10.1007/BF01162704}.
\bibitem{Rodemich} E. Rodemich, An inequality for univalent functions with
real coefficients, \emph{Notices Amer. Math. Soc.} \textbf{26} (January 1979),
abstract~763-30-11, A-68.
\href{https://web.archive.org/web/20220620050040id_/https://www.ams.org/journals/notices/197901/197901FullIssue.pdf}{Archived publisher scan}.
\bibitem{Jenkins1962} J.~A. Jenkins, The general coefficient theorem and
certain applications, \emph{Bull. Amer. Math. Soc.} \textbf{68} (1962), 1--9.
\href{https://doi.org/10.1090/S0002-9904-1962-10670-3}{doi:10.1090/S0002-9904-1962-10670-3}.
\bibitem{Jenkins1960} J.~A. Jenkins, On univalent functions with real
coefficients, \emph{Ann. of Math.} \textbf{71} (1960), 1--15.
\href{https://doi.org/10.2307/1969877}{doi:10.2307/1969877}.
\bibitem{AlePop} I.~A. Aleksandrov and G.~A. Popova, Extremal properties
of univalent holomorphic functions with real coefficients,
\emph{Sibirsk. Mat. Zh.} \textbf{14} (1973), 915--926;
English translation, \emph{Siberian Math. J.} \textbf{14} (1973), 635--643.
\href{https://doi.org/10.1007/BF00969900}{doi:10.1007/BF00969900}.
\end{thebibliography}
\end{document}
